%% kapa: exam=1 round=33 year=2022 period=1 form=A subject=economics subject_ko=경제학원론 seq=22 number=62
%% source: Q-Net 감정평가사 기출문제 (cst00309) artlSeq=5209997 file=raw/1008/5209997/2022년도 제33회 감정평가사 1차시험 1교시(민법,경제학,부동산학).pdf
%% conversion: kordoc 4.24.51 -> tools/gen1.py (auto); review=proofread-sonnet; flags=none
\begin{question}{62}
물가지수에 관한 설명으로 옳은 것을 모두 고른 것은?
\begin{framedtext}
ㄱ. GDP 디플레이터를 산정할 때에는 국내에서 생산되는 모든 최종 재화와 서비스를 대상으로 한다. \par
ㄴ. 소비자물가지수의 산정에 포함되는 재화와 서비스의 종류와 수량은 일정 기간 고정되어 있다. \par
ㄷ. 생산자물가지수를 산정할 때에는 기업이 생산 목적으로 구매하는 수입품은 제외한다.
\end{framedtext}
\begin{choices}
  \item ㄱ
  \item ㄴ
  \item ㄱ, ㄴ
  \item ㄱ, ㄷ
  \item ㄴ, ㄷ
\end{choices}
\end{question}
\answer{3}
